% =====================================================================
% Homework 03 --- Analise e otimizacao com Euler solver
% Incluido por main.tex. Tabelas em tex_eulerblock/ e figuras em
% resultados_eulerblock/ (mesmo nivel que este arquivo e o main.tex).
% Codigos em eulerblock/.
% =====================================================================

\section{Analise e Otimizacao com Euler Solver}
\label{sec:euler}

\subsection{Objetivo}

Explique nesta secao o papel do solver de Euler no laboratorio:
analise do perfil em escoamento inviscido, comparacao com o XFOIL,
estudo de malha, ou otimizacao do perfil.

\subsection{Modelo Numerico}

Descreva aqui:
\begin{itemize}
\item tipo de solver e ordem numerica;
\item malha adotada;
\item condicoes de contorno;
\item criterio de convergencia;
\item variaveis de projeto e restricoes, se houver otimizacao.
\end{itemize}

\begin{equation}
\min_{\mathbf{x}} \; f(\mathbf{x})
\qquad
\text{sujeito a} \qquad
\mathbf{g}(\mathbf{x}) \ge 0.
\end{equation}

\subsection{Resultados}

Inclua aqui os resultados principais desta etapa: historico de convergencia,
campo de pressao, comparacao de perfis, frente de iteracoes, ou historico
da otimizacao.

\begin{figure}[H]
\centering
% \includegraphics[width=0.85\textwidth]{resultados_eulerblock/exemplo.png}
\caption{Substitua esta legenda por uma descricao da principal figura da
etapa com o solver de Euler.}
\label{fig:euler-exemplo}
\end{figure}

\subsection{Discussao}

Comente a convergencia, o custo computacional, a sensibilidade a malha e a
comparacao entre os resultados do Euler solver e os obtidos no XFOIL.
